\chapter{室外空气环境与健康  【重点章节】}
\setcounter{chapter}{3}
\chapterintro{
本章为\keyword{重点考察章节}（6学时）。掌握大气层的结构及卫生学意义；掌握大气污染对健康的直接和间接危害（温室效应、臭氧层破坏、酸雨等）；掌握影响大气污染物浓度的因素；掌握主要大气污染物（SO$_2$、颗粒物、NOx、O$_3$等）的健康效应；理解AQI、大气卫生标准等。
}

% ----- 3.1 大气特征 -----
\section{大气的特征及卫生学意义}

\subsection{大气层结构}

\begin{table}[htbp]
\centering
\caption{大气圈五层结构及特征}
\begin{tabularx}{\textwidth}{l>{\raggedright\arraybackslash}p{2.5cm}>{\raggedright\arraybackslash}X}
\hline
\textbf{层次} & \textbf{高度范围} & \textbf{主要特征} \\
\hline
\imp{对流层} & 10--16 km & 气温水平/垂直运动；集中3/4大气质量；气象变化主要发生层；\textbf{与人类健康关系最密切} \\
\imp{平流层} & 16--55 km & 含\keyword{臭氧层}（20--30 km），吸收紫外线；同温层气温恒定 \\
中间层 & 55--85 km & 高度升高，气温下降 \\
热成层 & 85--800 km & 高度升高，气温上升 \\
逸散层 & $>$800 km & 大气稀薄，气体分子逸散 \\
\hline
\end{tabularx}
\end{table}

\begin{defbox}
\textbf{电离层（Ionosphere）}：约60--1000 km，覆盖中间层顶部、热成层和逸散层底部。由于太阳和宇宙射线的作用，大部分空气分子发生电离，具有较高密度的带电粒子，能将电磁波反射回地球，对全球无线电通讯有重大意义。
\end{defbox}

\subsection{大气的组成}

\begin{itemize}[leftmargin=*]
    \item \imp{干洁空气}：不含水汽和气溶胶
    \item \imp{水汽}
    \item \imp{气溶胶}：液体或固体微粒均匀分散在气体中形成的相对稳定的悬浮体系
\end{itemize}

\begin{defbox}
\textbf{雾与霾的区分：}
\begin{itemize}[leftmargin=*]
    \item \textbf{雾}：由大量悬浮在近地面空气中的微小水滴或冰晶组成，影响能见度，是水汽凝结（凝华）的产物。相对湿度 $>$ 90\%
    \item \textbf{霾}：由于空气中悬浮的大量颗粒物所致的水平能见度降低到10 km以下。相对湿度 $<$ 80\%
    \item 介于80\%--90\%之间为雾和霾共同作用的结果
\end{itemize}
\end{defbox}

\subsection{空气离子化}

\begin{defbox}
\textbf{空气离子化（Air Ionization）}：大气分子或原子形成带电荷的正负离子过程。\keyword{负离子}具有清洁空气、改善健康的积极作用。

\textbf{轻离子（Light Ion）}：带电气体分子，包括阳离子和阴离子。

\textbf{重离子（Heavy Ion）}：轻离子与空气中的悬浮颗粒或水滴结合形成。

清洁空气中轻离子浓度高，污染空气中重离子浓度高。空气中重离子数与轻离子数之比 $<$ 50 时，空气较清洁。
\end{defbox}

\subsection{太阳辐射的卫生学意义}
\begin{itemize}[leftmargin=*]
    \item \imp{紫外线（UV, 200--400 nm）}：
    \begin{itemize}
        \item UV-A（320--400 nm）：色素沉着作用
        \item UV-B（290--320 nm）：红斑作用、抗佝偻病作用
        \item UV-C（200--290 nm）：杀菌作用
        \item 过量可致皮肤癌
    \end{itemize}
    \item \imp{可见光（400--760 nm）}：综合作用于高级神经系统，提高视觉和代谢能力，平衡兴奋和镇静作用
    \item \imp{红外线（760--30000 nm）}：热效应；适量可促进新陈代谢和细胞增殖，过强可引起日射病和红外线白内障
\end{itemize}

\textbf{气象因素}（气温、气湿、气压和气流）对机体冷热感觉、体温调节、心血管功能、神经功能、免疫和新陈代谢功能有调节作用。

% ----- 3.2 大气污染 -----
\section{大气污染及污染物的转归}

\subsection{大气污染的来源}

\begin{enumerate}[leftmargin=*]
    \item \imp{工业污染}：燃料燃烧，排放SO$_2$、NOx、颗粒物等
    \item \imp{生活炉灶和采暖锅炉}：燃煤产生SO$_2$、烟尘
    \item \imp{交通运输}：汽车尾气（NOx、CO、VOCs、颗粒物）
    \item \imp{其他}：地面扬尘、意外事件等
\end{enumerate}

\subsection{影响大气污染物浓度的因素}

\begin{keybox}
\textbf{影响大气污染物浓度的三大因素：}

\textbf{一、污染源的排放情况}
\begin{enumerate}[leftmargin=*]
    \item \imp{排放量}：决定大气污染程度的最基本因素
    \item \imp{与污染源的距离}：
    \begin{itemize}
        \item \textbf{烟波着陆点}：烟气自烟囱排出后，向下风侧扩散稀释，接触地面的点。污染物浓度烟波着陆点最大，下风侧浓度随距离增大而降低
    \end{itemize}
    \item \imp{有效排出高度}：烟囱高度 + 烟气抬升高度。排出高度越高，污染物稀释程度越大。下风侧污染物最高浓度与有效排出高度的平方成反比
\end{enumerate}

\textbf{二、气象条件}
\begin{enumerate}[leftmargin=*]
    \item \imp{风和湍流}：水平输送和垂直扩散污染物；湍流强有利于稀释扩散

    \item \imp{风向频率图（风玫瑰图）}：按8或16个方位将一定时期的风向频率绘制成图。风向频率 = 某方向吹来风的次数 / 各方向吹来风的总次数 $\times$ 100\%。\textbf{主导风向}为风向频率最大的风向，其下风向地区污染最严重。只有一个主导风向时，居民区设在主导风向上风向，工业区设在下风向；有多个主导风向时，居民区设在最小风向的下风向，工业区设在其上风向。

    \item \imp{大气稳定度}：\\
    \begin{tabular}{ll}
    $\gamma > \gamma_d$（气温直减率 > 干绝热直减率）& 不稳定状态，波浪型 \\
    $\gamma = \gamma_d$ & 中性状态 \\
    $\gamma < \gamma_d$ & 稳定状态，锥型 \\
    $\gamma < 0$ & \keyword{逆温}，极稳定状态，扇型 \\
    \end{tabular}

    \item \imp{逆温（Temperature Inversion）}：大气温度随高度升高而上升，上层温度高于下层，极不利于污染物扩散。常见类型：
    \begin{itemize}
        \item \textbf{辐射逆温}：无风少云的夜晚，地面通过辐射失去热量而冷却，近地空气随之冷却，上层降温较慢
        \item \textbf{下沉逆温}：空气压缩增温，上层空气下沉落入高气压团中受压变热
        \item \textbf{地形逆温}：局部地理条件形成（如盆地和山谷），冷空气沿山坡聚集在谷底，暖空气受挤压抬升
    \end{itemize}

    \item \imp{气压、气湿、降水}：降水能清除部分污染物；气湿大加重局部污染，利于二次污染物形成（如SO$_2$溶于水易被氧化为硫酸，促进酸雨形成）
\end{enumerate}

\textbf{三、地形条件}
\begin{enumerate}[leftmargin=*]
    \item 山地和谷地：容易形成逆温，不利扩散
    \item \imp{海陆风}：白天升温陆地$>$水面，形成海风；夜晚散热陆地$>$水面，形成陆风。若污染源在岸边，白天可能污染岸上居民区
    \item \imp{城市热岛效应}：人口密集的城市热量散发远大于郊区，城市热空气上升，郊区冷空气补充，把郊区排放的污染物引入城市，造成城市大气污染
\end{enumerate}
\end{keybox}

\subsection{大气污染物的种类}

\textbf{按属性分：}物理性（噪声、电离辐射、电磁辐射等）、化学性和生物性（经空气传播的病原微生物和植物花粉等）三类，其中以化学性污染种类最多、污染范围最广。

\textbf{按存在状态分：}气态和气溶胶。气态污染物包括气体和蒸汽；气溶胶中分散的各种颗粒即大气颗粒物。

\subsection{大气物理性污染与生物性污染}

\begin{defbox}
\textbf{物理性污染：}
\begin{itemize}[leftmargin=*]
    \item \imp{噪声污染}：影响听觉系统（耳鸣、听阈位移、高频听力丧失、耳聋）、神经系统、心血管系统
    \item \imp{光污染}：白亮污染、人工白昼、彩光污染，影响视觉系统、生物钟
    \item \imp{放射性辐射污染}：肺癌（氡）
    \item \imp{电磁辐射污染}
\end{itemize}

\textbf{生物性污染：}
\begin{itemize}[leftmargin=*]
    \item 微生物：呼吸道传染疾病（集中在冬春季）
    \item 变应原：植物花粉、真菌孢子等，引起哮喘、过敏
    \item 其他生物质：柳絮、梧桐叶等生物性尘埃，引起过敏、哮喘、鼻炎、皮肤瘙痒
\end{itemize}
\end{defbox}

\subsection{二次污染}

\begin{defbox}
\textbf{二次污染（Secondary Pollution）}：某些污染物进入大气环境后，本身没有消失，一旦转移后环境发生变化，该污染物可回到原先环境造成二次污染。

\textbf{二次污染物}：由一次污染物在大气中经化学或光化学作用转变成的毒性更大的化学物质。如SO$_2$ → 硫酸雾，NOx + VOCs + 日光 → 光化学烟雾（O$_3$、PAN）。

\textbf{大气污染物的转归：}①自净——扩散和沉降、氧化和中和、植物吸附；②转移——下风侧转移、向其他环境介质转移、向平流层转移；③转化——形成二次污染和二次污染物。
\end{defbox}

% ----- 3.3 大气污染的健康影响 -----
\section{室外空气污染对健康的影响}

\subsection{直接危害}

\begin{keybox}
\textbf{（一）急性危害}

1. \imp{烟雾事件}：

\begin{table}[H]
\centering
\caption{煤烟型烟雾与光化学烟雾的比较}
\begin{tabularx}{\textwidth}{lXX}
\hline
\textbf{比较要点} & \textbf{煤烟型烟雾（London Smog）} & \textbf{光化学烟雾（Photochemical Smog）} \\
\hline
主要污染物 & SO$_2$、颗粒物 & O$_3$、PAN、醛类（光化学氧化剂） \\
污染来源 & 燃煤 & 汽车尾气（NOx + VOCs） \\
气象条件 & 冬季、早晨、低温、辐射逆温、高湿有雾 & 夏秋季、午后、高温、下沉逆温、低湿、紫外线强烈 \\
地形 & 山谷或盆地 & 三面环山，西面临海（如洛杉矶） \\
主要健康效应 & 刺激呼吸道 & 刺激眼睛和呼吸道 \\
\hline
\end{tabularx}
\end{table}

2. \imp{生产事故}：如印度博帕尔异氰酸甲酯泄漏事件（1984）。

\textbf{（二）慢性危害}
\begin{enumerate}[leftmargin=*]
    \item 影响呼吸系统功能：反复刺激→咽炎、喉炎、气管炎→气道狭窄、纤毛损伤、肺泡扩张→COPD；SO$_2$、O$_3$、NOx吸附于颗粒物成为变态反应原，引起支气管哮喘
    \item 心血管疾病：动脉粥样硬化、缺血性心脏病、心律失常、心衰、心搏骤停
    \item 降低免疫功能
    \item 引起变态反应
    \item 致癌作用（如多环芳烃与肺癌）
\end{enumerate}
\end{keybox}

\subsection{间接危害}

\begin{keybox}
\textbf{1. 温室效应（Greenhouse Effect）}
\begin{itemize}[leftmargin=*]
    \item 原因：CO$_2$、CH$_4$等温室气体吸收地表热辐射，使大气增温
    \item 危害：海平面上升；促进酸雨形成；热相关疾病增加；虫媒传染病流行范围扩大
\end{itemize}

\textbf{2. 臭氧层破坏}
\begin{itemize}[leftmargin=*]
    \item 原因：CFCs（氯氟烃）、N$_2$O、CCl$_4$、哈龙等消耗臭氧层物质
    \item 危害：皮肤癌、白内障发生率增加；影响生态系统
\end{itemize}

\textbf{3. 酸雨}
\begin{itemize}[leftmargin=*]
    \item 定义：降水pH值 $<$ 5.6即为酸雨
    \item 我国以\keyword{硫酸型酸雨}为主，主要因燃煤排放SO$_2$和NOx
    \item 危害：土壤和植物损害；影响水生生态系统；腐蚀建筑物和文物；影响人类健康
    \item 防治：控制SO$_2$和NOx排放
\end{itemize}

\textbf{4. 大气棕色云团（ABC）}
\begin{itemize}[leftmargin=*]
    \item 区域性大气污染现象，由颗粒物和气体污染物共同形成
\end{itemize}
\end{keybox}

\subsection{主要大气污染物及其健康影响}

\begin{table}[htbp]
\centering
\caption{主要大气污染物及其健康效应}
\begin{tabularx}{\textwidth}{lX}
\hline
\textbf{污染物} & \textbf{主要健康效应} \\
\hline
\imp{颗粒物（PM）} & 我国多数城市的\keyword{首要污染物}。呼吸系统：堵塞、刺激肺泡壁、氧化应激炎症反应；心血管：干扰自主神经→心律不齐、诱发血栓、刺激呼吸道释放细胞因子损伤血管。粒径$>$5$\mu$m沉积上呼吸道，$<$5$\mu$m进入细支气管和肺泡 \\
\imp{SO$_2$} & 阻滞于上呼吸道→亚硫酸盐/硫酸氢盐；气管和支气管收缩，气道阻力↑；吸附SO$_2$的颗粒物是变态反应原→支气管哮喘；促癌（增加BaP致癌性） \\
\imp{NOx} & NO$_2$难溶于水，主要作用于深部呼吸道、细支气管和肺泡；形成亚硝酸盐和硝酸盐→高铁血红蛋白→组织缺氧；与SO$_2$、O$_3$有相加和协同作用 \\
\imp{O$_3$} & 二次污染物（占光化学氧化剂90\%以上），光化学烟雾的指示物；体内→超氧阴离子→呼吸道上皮脂质过氧化→炎症；激活NF-$\kappa$B \\
\imp{CO} & 与Hb亲和力比O$_2$大200--300倍→HbCO→降低血液携氧能力→阻碍氧释放→组织缺氧；还可与肌红蛋白、细胞色素氧化酶及P450结合；急性中毒以神经系统症状为主；\textbf{胎儿对CO毒性比成人敏感}——孕妇吸烟可致低体重儿、围产期死亡增高及婴幼儿神经行为障碍 \\
\imp{Pb（铅）} & 全身性毒物，90\%贮存于骨骼。神经系统（神经衰弱→周围神经炎）；造血系统（抑制Hb合成→贫血）；可透过胎盘屏障影响胚胎发育。\textbf{儿童为敏感人群}（户外活动多、摄入量大、吸收率高、血脑屏障发育不完全）——铅选择性蓄积于脑海马部位，损害神经细胞形态和功能，表现为注意力不集中、记忆力降低、学习能力下降等。交警为铅职业接触高危人群 \\
\imp{PAHs（多环芳烃）} & 致癌（如苯并[a]芘与肺癌） \\
\imp{二噁英} & 主要来源于垃圾焚烧、工业生产；吸附于颗粒物→沉降到水体和土壤→通过食物链富集进入人体。食物是人体内二噁英的主要来源 \\
\hline
\end{tabularx}
\end{table}

\begin{defbox}
\textbf{空气动力学等效直径（Dp）}：在气流中，被测颗粒的空气动力学效应与密度为1的球形标准颗粒物相同时，该标准颗粒物粒径即被测颗粒物的Dp。可表达颗粒物在空气中的停留时间、沉降速度、进入呼吸道的可能性及在呼吸道的沉降部位。

\textbf{TSP（总悬浮颗粒物）}：粒径 $\leq$ 100 $\mu$m的颗粒物，包括液体、固体或液体与固体结合并悬浮在空气介质中的颗粒。

\textbf{IP（可吸入颗粒物）/ PM$_{10}$}：粒径 $\leq$ 10 $\mu$m，能进入人体呼吸道，长期漂浮在空气中。

\textbf{PM$_{2.5}$（细颗粒物）}：粒径 $\leq$ 2.5 $\mu$m，可进入肺泡和血液循环，易吸附有毒有机物和重金属元素，对健康危害极大。

\textbf{PM$_{0.1}$（超细颗粒物）}：粒径 $\leq$ 0.1 $\mu$m。
\end{defbox}

% ----- 3.4 调查与监测 -----
\section{大气污染和气候变化对健康影响的调查与监测}

\subsection{环境空气质量功能分区与标准分级}

\begin{keybox}
\textbf{环境空气质量功能分区：}
\begin{itemize}[leftmargin=*]
    \item \imp{一类区}：自然保护区、风景名胜区和其他需要特殊保护的地区 → 执行\imp{一级标准}（长期接触不发生任何危害影响）
    \item \imp{二类区}：居住区、商业交通居民混合区、文化区、工业区和农村地区 → 执行\imp{二级标准}（长期和短期接触不发生伤害）
\end{itemize}

\textbf{污染物浓度限值类型：}
\begin{itemize}[leftmargin=*]
    \item \imp{1小时平均浓度}：任何1h内平均浓度的最高容许值。针对能在短期内产生刺激、过敏或中毒等急性危害的物质
    \item \imp{24小时平均浓度（日平均）}：任何自然日24h平均浓度的最高容许值
    \item \imp{年平均浓度}：任何日历年内各日平均浓度的算术均值的最高容许值。针对有慢性作用的物质，防止慢性和潜在危害
\end{itemize}
\end{keybox}

调查内容：①污染源调查；②污染状况监测；③人群健康调查。

\subsection{采样布点方法}

\begin{table}[htbp]
\centering
\caption{大气污染监测采样布点方法}
\begin{tabularx}{\textwidth}{lX}
\hline
\textbf{污染源类型} & \textbf{布点方法} \\
\hline
\imp{点源污染} & ①四周布点：以污染源为中心，8个方位不同距离同心圆布点；②扇形布点：主导方向下风侧划3--5个方位布置；③捕捉烟波布点：随烟波方向变动布点（不固定） \\
\imp{面源污染} & ①按城市功能分区布点（每种类型2--3个采样点+清洁对照点）；②几何状布点（方形或三角形网格）；③综合因素布点 \\
\imp{线源污染} & 采样口离地面高度2--5 m，距道路边缘$\leq$20 m \\
\hline
\end{tabularx}
\end{table}

% ----- 3.5 AQI -----
\section{大气质量标准}

\begin{defbox}
\textbf{日平均最高容许浓度}：任何一天内多次测定的平均浓度的最高容许值。

\textbf{一次最高容许浓度}：任何一次短时间采样测量结果的最高容许值。
\end{defbox}

\begin{tipbox}
\textbf{AQI（空气质量指数）}：属于分段线性函数型大气质量指数。参数包括SO$_2$、NO$_2$、PM$_{10}$、PM$_{2.5}$、CO、O$_3$共6项。

\textbf{空气质量分指数（IAQI）}：各污染物的分指数。AQI取各IAQI中的最大值，该污染物即为\imp{首要污染物}。

\textbf{确定方法：}
\begin{itemize}[leftmargin=*]
    \item AQI $>$ 50时，IAQI最大的污染物为首要污染物（若多项并列则并列为首要污染物）
    \item IAQI $>$ 100的污染物为\imp{超标污染物}
\end{itemize}

\textbf{优缺点：}优点——突出当日单个污染物的作用，便于发现首要污染物及其对健康的影响；缺点——对各污染物指数无综合，仅给出最大分指数，不能反映其他污染物的综合影响。
\end{tipbox}

\section{气候变化与极端天气的健康影响}

\subsection{气候变化（Climate Change）}

\begin{keybox}
\textbf{气候变化对健康的直接影响：}
\begin{itemize}[leftmargin=*]
    \item \imp{高温热浪}：可改变人体生理功能，增加中暑、热射病、急性肾衰竭等急性病风险
    \item \imp{洪涝与风暴潮}：造成溺水和创伤，导致死亡率和伤残率增加
    \item \imp{野火燃烧和沙尘天气}：通过携带大气颗粒物等污染物直接影响健康
\end{itemize}

\textbf{气候变化对健康的间接影响：}
\begin{enumerate}[leftmargin=*]
    \item \imp{传染病风险加剧}：温度升高与降水模式改变使昆虫传播媒介和病原体的地理分布范围扩大、繁殖季节延长，导致病毒传播到纬度与海拔更高的地区
    \item \imp{粮食安全}：干旱、洪涝、高温热浪导致农作物产量减少及粮食营养成分下降
    \item \imp{水资源短缺}：影响饮用水供应与水质，导致卫生条件恶化，引发肠道传染病
    \item \imp{过敏性疾患增加}：气候变暖使真菌孢子、花粉等浓度增高
\end{enumerate}
\end{keybox}

\subsection{极端天气}

\begin{defbox}
\textbf{极端天气事件}：显著偏离正常气候状况，对人类社会和自然生态系统构成严重威胁的高强度、低频次天气现象。特征为强度高、影响范围广、持续时间长。

\begin{itemize}[leftmargin=*]
    \item \imp{高温热浪}：增加中暑、热射病、热痉挛等热相关疾病风险
    \item \imp{寒潮}：增加心血管疾病发作风险，削弱皮肤和呼吸道黏膜免疫屏障
    \item \imp{极端强降水}：引发山洪、内涝，造成人员伤亡；污染清洁水源，导致病原微生物和媒介生物快速孳生，造成传染病暴发流行
    \item \imp{洪涝}：引发心理疾病（抑郁、焦虑）；水源污染和细菌孳生（霍乱、痢疾等）引发肠道和呼吸道感染；破坏公共设施造成救护延误，增加慢性非传染性疾病死亡风险
    \item \imp{干旱}：粮食减产，受灾人群营养不良、营养缺乏风险增加
    \item \imp{野火}：烟雾中有害颗粒物和气体引发或加剧哮喘、COPD等呼吸系统疾病，增加心血管疾病风险
    \item \imp{沙尘暴}：加重呼吸系统疾病，引起眼部刺激和结膜炎，增加过敏反应
    \item \imp{台风}：机体损伤（软组织挫伤、颅脑损伤、骨折），增加介水传染病和食源性传染病风险
\end{itemize}
\end{defbox}

\section{大气卫生防护措施}

\begin{itemize}[leftmargin=*]
    \item \textbf{规划措施}：合理安排工业布局和城镇功能分区；加强绿化；加强局部污染源管理
    \item \textbf{工艺措施}：控制燃煤污染（清洁能源、集中供热）；加强废气治理（除尘、脱硫、脱硝）
\end{itemize}

% ----- 3.7 复习题 -----
\section{复习题}

\begin{enumerate}[leftmargin=*, itemsep=8pt]
    \item \textbf{【名词解释】}光化学烟雾；逆温；酸雨
    \item \textbf{【名词解释】}TSP；PM$_{10}$；PM$_{2.5}$
    \item \textbf{【名词解释】}一次污染物与二次污染物（大气中）
    \item \textbf{【名词解释】}温室效应
    \item \textbf{【简答题】}试述大气稳定度对大气污染扩散的影响。
    \item \textbf{【简答题】}请简述影响大气中污染物浓度的因素。
    \item \textbf{【简答题】}大气污染对健康的直接危害有哪些？间接危害有哪些？
    \item \textbf{【简答题】}臭氧层破坏的原因及其健康危害是什么？
    \item \textbf{【简答题】}简述光化学烟雾与煤烟型烟雾的主要区别。
    \item \textbf{【简答题】}简述酸雨的形成原因、危害及防治对策。
    \item \textbf{【简答题】}简述大气物理性污染和生物性污染的主要类型及健康影响。
    \item \textbf{【简答题】}简述气候变化对健康的直接和间接影响。
    \item \textbf{【简答题】}简述极端天气事件的类型及其健康危害。
    \item \textbf{【简答题】}简述大气卫生防护的主要措施。
    \item \textbf{【非标准答案题】}结合我国当前大气污染现状，谈谈你对大气污染防治与人群健康保护的思考。
\end{enumerate}

\subsection*{参考答案要点}

\begin{enumerate}[leftmargin=*, itemsep=5pt]
    \item \textbf{光化学烟雾}：汽车尾气中NOx和VOCs在日光紫外线照射下经光化学反应生成的浅蓝色烟雾，主要含O$_3$、PAN等强氧化性物质。\textbf{逆温}：大气温度随高度升高而上升，上层气温高于下层，极不利于污染物扩散。\textbf{酸雨}：降水pH值$<$5.6，主要由SO$_2$和NOx转化形成。

    \item \textbf{TSP}：总悬浮颗粒物，粒径$\leq$100$\mu$m。\textbf{PM$_{10}$}：可吸入颗粒物，粒径$\leq$10$\mu$m，能进入呼吸道。\textbf{PM$_{2.5}$}：细颗粒物，粒径$\leq$2.5$\mu$m，可进入肺泡和血液循环。

    \item （参照第1章定义）

    \item 大气中CO$_2$、CH$_4$等温室气体吸收地表热辐射，使大气增温的作用。

    \item （答题要点：$\gamma$与$\gamma_d$的比较→不稳定/中性/稳定状态；逆温条件下的极稳定状态；各种稳定度对应的烟型及其对污染物扩散的影响）

    \item （答题要点：三大因素——污染源排放情况（排放量、高度、距离）、气象条件（风、湍流、大气稳定度、逆温、降水）、地形条件（山地谷地、海陆风、热岛效应））

    \item 直接危害：急性（烟雾事件、生产事故）和慢性（呼吸系统疾病、心血管疾病、免疫功能降低、致癌）。间接危害：温室效应、臭氧层破坏、酸雨、大气棕色云团。

    \item 原因：CFCs、N$_2$O、CCl$_4$等消耗臭氧层物质释放。危害：皮肤癌和白内障增多，影响生态系统。

    \item 光化学烟雾：汽车尾气+日光，含O$_3$、PAN，夏秋季午后多，刺激眼睛。煤烟型：燃煤+逆温，含SO$_2$、烟尘，冬季清晨多，刺激呼吸道。

    \item 成因：燃煤和石油燃烧排放SO$_2$和NOx，在大气中氧化转化为硫酸和硝酸。危害：土壤酸化、水生生态破坏、建筑物腐蚀、健康影响。防治：控制SO$_2$和NOx排放。

    \item 规划措施（合理布局、绿化、局部源管理）和工艺措施（控制燃煤、加强废气治理）。

    \item （开放性问题，可围绕清洁能源、产业结构调整、多部门协作、公众参与等方面展开。）
\end{enumerate}

\newpage

% =====================================================================
%  CHAPTER 4: 室内环境与健康
% =====================================================================
\newpage